\section{A tail-content restriction for the remaining long ordinary shapes}
\label{sec:ordinary-alternating-long-row}

\paragraph{Status.}
New root proof, read in full by its author and independently by
/root/g34\_eightfolds; both whole-proof reviews PASS on 21 November
2026. Separate exact arithmetic certificates also passed. This is one further uniform shape reduction;
no claim that it settles all long shapes is made.

\begin{lemma}[HD-ALT-05a: exact tail-content estimates]
\label{HD-ALT-05a}
Let $n\ge15$ and let $\lambda=(n-k,\mu)$, where $\mu\vdash k$
and $0<k<n/4$. Put $t=\sum(j-i)$ and $S=\sum(j-i)^2$ over
the boxes of $\lambda$. Then
\begin{align}
 S&\le\sum_{j=1}^{n-k-1}j^2+\sum_{j=1}^k j^2,               \tag{ALTL.1}\\
 0<t&\le\binom n2-k(n-k+1).                                \tag{ALTL.2}
\end{align}
For the three-cycle ratio of its irreducible alternating restriction,
\[
 \frac{1-u}{3}\ge\frac{k(n-1-k)}{(n-1)(n-2)}.
                                                               \tag{ALTL.3}
\]
If also $k\ge n/7$, then its double-transposition ratio satisfies
\[
                         v<7/12.
                                                               \tag{ALTL.4}
\]
Every nonscalar finite-order linear normalizer in that range consequently
satisfies $\nu(T)\ge5a/48$, where $a=\dim S^\lambda$.
\end{lemma}

\begin{proof}
Choose a standard tableau of the tail diagram $\mu$, which orders its
boxes by successive valid Young-diagram additions. If its $r$th box
is $(i,j)$ in tail coordinates, then $i,j\le r$. Its content in
$\lambda$ is $j-i-1$, lying between $-r$ and $r-2$. Its square is
therefore at most $r^2$. Adding these bounds and the exact first-row
squares proves (ALTL.1). The same lower bound for each tail content
gives
\[
 t\ge\binom{n-k}{2}-\frac{k(k+1)}2
     =\frac{n(n-2k-1)}2>0.
\]
The tail content sum is its own ordinary content sum minus $k$.
The ordinary sum for a partition of $k$ is at most $\binom k2$:
write it as the sum of row binomials minus the sum of column
binomials, and use $\sum_i\binom{\mu_i}{2}\le\binom k2$.
Hence
\[
 t\le\binom{n-k}{2}+\binom k2-k
       =\binom n2-k(n-k+1),
\]
proving (ALTL.2). Substituting (ALTL.1) into the exact content
formula (ALTQ.3), then simplifying, gives (ALTL.3).

Now suppose $n/7\le k<n/4$. On this interval both
$k(n-1-k)$ and $k(n+1-k)$ increase with $k$. Thus
\[
 \frac{1-u}{3}\ge\frac{n(6n-7)}{49(n-1)(n-2)}
                 >\frac6{49},\qquad
 0<t\le\frac{n(37n-63)}{98}.
                                                               \tag{ALTL.5}
\]
The first strict inequality follows on subtraction because the
numerator difference is $11n-12>0$. Cauchy's inequality as in
(ALTQ.6) now gives
\[
 v\le\frac{(37n-63)^2}{2401(n-1)(n-2)}
                +\frac4{(n-2)(n-3)}<\frac7{12}.
                                                               \tag{ALTL.6}
\]
After multiplication by the positive denominator
$12\cdot2401(n-1)(n-2)(n-3)$, the final strict inequality
is equivalent to positivity of
\[
 P(n)=379n^3+4386n^2-145831n+157290.
\]
Writing $n=15+x$ gives
\[
 P(15+x)=379x^3+21441x^2+241574x+235800>0
                         \quad(x\ge0).
\]
This proves (ALTL.4) for every $n\ge15$. Since
$6/49>5/48$, the complete normalizer estimate HD-ALT-03a gives
$\nu(T)\ge a\min\{(1-u)/3,(1-v)/4\}\ge5a/48$.
The long first row prevents a self-conjugate label, so its
alternating restriction is irreducible and has degree $a$.
\end{proof}

\begin{theorem}[HD-ALT-05b: the remaining row or column exceeds six sevenths]
\label{HD-ALT-05b}
In the transitive ordinary alternating-core setting of HD-ALT-04b,
every factor label satisfies
\[
                    \max(\lambda_1,\lambda'_1)>6n/7.
                                                               \tag{ALTL.7}
\]
The same allowance for arbitrary compatible outers, scalar phases,
tensor permutations and larger inducing character fields applies.
\end{theorem}

\begin{proof}
By HD-ALT-04b any remaining label has a row or column longer
than $3n/4$. Transposing if necessary, write it as $(n-k,\mu)$,
with $0<k<n/4$. A nontrivial ordinary factor has $k>0$.
If (ALTL.7) failed, then $k\ge n/7$. Lemma HD-ALT-05a gives
the bound $5a/48$ on each nonscalar factor normalizer. It passes
to factor-preserving tensor operators, while every nontrivial
factor permutation has the stronger quarter bound. Genuine
nonidentity scalar inertia moves the whole space.

The real-factor and nontrivial-multiplier argument of HD-ALT-04b
still applies. At least two of the at least twelve genuine
inertias have ordinary moving rank at least $\lceil e/2\rceil$.
The exact total-moving-rank identity therefore implies
\[
 2\lceil e/2\rceil+10\lceil5e/48\rceil\le2e+6,
 \qquad \frac{e}{24}\le6,
 \qquad e\le144.
                                                               \tag{ALTL.8}
\]
Every factor here has degree at least $2n^3/27\ge250$.
Use the same published precise low-degree classification consequence
\cite[proof of Theorem~9.7, p.~312]{HD:GKT2019} as in HD-ALT-04.
All three exceptional shapes have a row or column of length
at least $n-2>6n/7$ for $n\ge15$, and so cannot occur.
There is no factor-of-two loss on restriction now: the row
longer than $3n/4$ ensures the partition is not self-conjugate.
Thus $e\ge250$, contradicting (ALTL.8).
\end{proof}

\paragraph{Strictness and geometric meaning.}
At $n=15$ the shape $(12,3)$ lies outside the three-quarter
region and inside the newly excluded range, so this is a strict
further restriction. The result again concerns actual finite
local systems and the factorization of the Weil period map.
It narrows the remaining supplier gate for condition A, cote~16
scan~050, without producing a detecting cycle or altering the
condition I/projector distinctions in scan~049, cote~44 Theorem~4,
and scans~048--049. No new assertion about motivic $L$-values
or weights is inferred. The still longer shapes remain open
except where the separate standard, exterior-square and
unordered-pair exclusions apply.
